\section{Focus at the resolution: the coefficient stays one}
\label{sec:res}

\begin{proposition}[line form] \proofst
\label{prop:Jres}
Let $\ell\ge55.26$ and let the ripple satisfy $\eta\le\bar\eta=\tfrac14$.
Let a state satisfy $\tfrac12<b\le1$, $\mu>0$, $\sigma^3\le r\le\tfrac14$, $\kappa>0$, and let $P>0$ satisfy (A), (K), (CM) of the brightness lemma of~\cite{I}.
Then for every $a>0$ and every $\gamma\in[0,0.49]$ the lines of the option Fz satisfy $\Lambda_{\rm Fz}(a)\ge0$, with the target $\gamma\sigma\,\mathrm{den}$ (coefficient one).
\end{proposition}

\begin{proof}[Proof of Theorem~\ref{thm:Jres}]
If $G>0$, the lines of option Fz (valid for $b\in(\tfrac12,1]$, Sec.~V of~\cite{I}) and Proposition~\ref{prop:Jres} with $\gamma=0.49$ give $(a)+(a')\ge0.49$.
If $G=0$, Lemma~$G=0$ of~\cite{I} gives $(a)+(a')\ge0.1995$; it concerns a window without a ripple in the focus and uses no bound on $\eta$.
Both exceed $\Gamma_\ell\le0.1462$.
\end{proof}

\emph{Why the coefficient stays one.}
Option Fz uses only the two background lines P1, P4z and the ripple-free part of P3z; the gain P2z enters only to absorb $\mathrm{d4}+\mathrm{N4G}$.
No gain carries the sandwich factor $\lambda$, so nothing has to be scaled; the ripple enters only through the remainders, which we bound uniformly at $\bar\eta=\tfrac14$, and through $k_2=1/(2(1-\eta))\le\tfrac23$ in the target.

\subsection{Proof of Proposition~\ref{prop:Jres}}
The steps are those of Sec.~V of~\cite{I}, with $\eta=\tfrac14$ in every remainder.
Every term below is a product of non-negative factors that do not increase with $\ell$ ($1/\ell$, $\sigma$, $\ell\sigma$, $\sqrt\sigma$, $U_1^2/\ell$, $X_1^2/\ell$; $\eta=\tfrac14$ is fixed), so $\ell=55.26$ is the worst case.
Every number is an exact rational upper or lower bound: $e\ge2.71828$, $e^{-x}\le2.71828^{-\lfloor x\rfloor}/\sum_{j\le3}f^j/j!$ ($f=x-\lfloor x\rfloor$), $\sqrt{0.2}\le0.44722$, $\sqrt\pi\ge1.7724$, $\sqrt2\ge1.4142$, $P_1\le0.8829$ ((F1$'$) of~\cite{I}, which does not involve $\eta$); each square root is checked by squaring.
\begin{itemize}
\item[(R1$'$)] N1 $=\varepsilon_1(r)$P1 with $\varepsilon_1\le\bar\eta k_E\cdot1.8\le0.2013$; the other L-lines are at most $e_L$P1, $e_L\le2.06\times10^{-14}$ (the largest terms are A2b and B2b, of order $\sqrt\sigma\,\bar\eta^2$); $\mathrm{d4}+\mathrm{N4G}\le1.4\times10^{-4}\,\mathrm{es}_\sigma\,$P2z.
\item[(R2$'$)] hneg1 $=\varepsilon_h(r)$P4z, $\varepsilon_h\le2k_E(\bar\eta+0.184)/\omega\le0.3882$.
\item[(R3$'$)] The S-lines are at most $e_S$P4z ($a\le1$) or $e_S$P1 ($a\ge1$), $e_S\le4.0\times10^{-13}$ ($c_S\le0.1985$). So $e_L+e_S\le e_2=10^{-12}$, with room $5.8\times10^{-13}$.
\item[(R4$'$)] $\mathrm{N3}+\mathrm{M0Xi}+\mathrm{A1v}+\mathrm{pair}\le1.9165\,\sigma\le2\sigma$. (Paper~I had $1.86\sigma$; at $\bar\eta=\tfrac14$ this bound fails, and $2\sigma$ is still covered by the tail term below for $\ell\ge1.5$.)
\item[(R5$'$)] $\mathrm{FzK}+\mathrm{A2k}+\mathrm{B2a}\le0.0682\,\ell\sigma\le\ell\sigma$.
\end{itemize}
The chain $\varepsilon_1$, $\varepsilon_h$, $\theta^2$, $k_1$ and the two margins below are checked step by step in exact rationals, each from the printed bounds before it (\code{dy\_room2\_hand.py}).
Hence, as in~\cite{I}, for every $a>0$,
\begin{equation}
\Lambda_{\rm Fz}(a)\ \ge\ \theta\sqrt b-\gamma k_1b-\gamma k_2G-(2\ell+\tfrac12)\sigma,
\end{equation}
with $\theta^2\ge(1-\varepsilon_1-e_2)(1-\varepsilon_h-e_2)\omega\ge0.48864$, $k_1=1+(\bar\eta+0.184)k_E\le1.19413$, $k_2=\tfrac23$ and $G\le P_1/3$.
The right side is concave in $b$; at the ends, in squared form with $\gamma=0.49$,
\begin{gather}
b=1:\quad \theta^2-\bigl(\gamma k_1+\gamma k_2P_1/3+\text{tail}\bigr)^2\ \ge\ 0.0245,\\
b=\tfrac12:\quad \theta^2-2\bigl(\gamma k_1/2+\gamma k_2P_1/3+\text{tail}\bigr)^2\ \ge\ 0.186 .
\end{gather}
$\Lambda_{\rm Fz}$ decreases in $\gamma$, so the bound holds on $[0,0.49]$.
Case~I of the rule holds throughout: $c_1\sqrt{1/2}/\Gamma_2-P_1/3\ge7.15$ with $c_1=0.77$, $\Gamma_2=\Gamma_\ell/2$.
\qed

\emph{What is hand and what is computed.}
Every inequality above is closed by hand: the arithmetic of the rational bounds is done exactly (by \code{fractions}), and nothing is evaluated in floating point or interval arithmetic.
The same quantities were also enclosed by interval arithmetic on $[55.26,\infty)$ as one piece, and the hand bounds lie above the interval values.
At $\gamma=\Gamma_{55.26}\le0.1462$ the margin at $b=1$ is $0.496$, so the law with coefficient one has room to spare.
The proof gives $0.49$ in place of the $0.6$ of~\cite{I} for a weak ripple; for $\eta\le1/\ell$ the $0.6$ of~\cite{I} stands.
With $\gamma=0.51$, or with $\bar\eta=0.3$, the bound at $b=1$ is negative.

\emph{In the observer's words.}
A focus as narrow as the resolution sees less than a hundredth of a period.
The ripple can only darken what it sees.
Even the strongest ripple of the class leaves the two background terms: wherever the ripple reaches the focus ($G>0$) they keep $(a)+(a')$ at $0.49$ or more, more than three times the $\Gamma_\ell\le0.146$ that the law asks for; where it does not, $(a)+(a')\ge0.1995$, still above it.
